\documentclass{article}
\usepackage[margin=2cm]{geometry}
\usepackage{amsmath}
\usepackage{graphicx}
\usepackage{booktabs}
\usepackage[utf8]{inputenc}
\usepackage{ragged2e}
\setlength{\emergencystretch}{3em}
\begin{document}
\sloppy

\subsection*{0.1Mass deformations}


In the following, we consider the coset with n = 1, i.e. a single vector multiplet. The coset
representative is expressed in terms of four scalars \$', i = 0, 1, 2, 3 via


\begin{equation}
L = \prod _ { i = 0 } ^ { 3 } e ^ { \phi ^ { i } K ^ { * } } \qquad ( 1 )
\end{equation}


SU(2)R singlet, while the other three scalars \$r form an SU(2)R triplet. The scalar potential
for this specific case can be obtained from and takes the following form


\begin{equation}
\begin{array} { r } { V ( \sigma, \phi ^ { i } ) = - g ^ { 2 } e ^ { 2 \sigma } + \frac { 1 } { 8 } m e ^ { - 6 \sigma } \bigg [ - 3 2 g e ^ { 4 \sigma } \cosh \phi ^ { 9 } \cosh \phi ^ { 1 } \cosh \phi ^ { 2 } \cosh \phi ^ { 3 } + 8 m \cosh ^ { 2 } \phi ^ { 0 } } \\ { \qquad }\end{array}
\end{equation}


The supersymmetric AdS6 vacuum is given by setting g = 3m and setting all scalars to
vanish. The masses of the linearized scalar fluctuation around the AdS vacuum determine
the dimensions of the dual scalar operators in the SCFT via


\begin{equation}
m ^ { 2 } l ^ { 2 } = \Delta ( \Delta - 5 ) \qquad ( 3 )
\end{equation}


where l is the curvature radius of the AdS6 vacuum. For the scalars at hand, one finds.


\begin{equation}
m _ { \sigma } ^ { 2 } l ^ { 2 } = - 6 ~ m _ { \phi } ^ { 2 } l ^ { 2 } = - 4 ~ m _ { \phi ^ { \prime } } ^ { 2 } l ^ { 2 } = - 6 ~, ~ ~ r = 1, 2, 3 ~
\end{equation}


Hence the dimensions of the dual operators are


\begin{equation}
\Delta _ { \mathcal { O } _ { \varphi } } = 3, ~
\end{equation}


In these CFT operators were expressed in terms of free hypermultiplets (i.e. the singleton
sector). The case of n = 1 corresponds to having a single free hypermultiplet, consisting of.
four real scalars qA and two symplectic Majorana spinors l. Here I = 1, 2 is the SU(2)R
R-symmetry index and A = 1,2 is the SU(2) flavor symmetry index. The gauge invariant.
operators appearing in are related to these fundamental fields as follows,


\begin{equation}
\begin{array} { r l } { \mathcal { O } _ { \sigma } = ( q ^ { * } ) ^ { A } q _ { \mathrm { ~ A } } ^ { I }, \qquad \mathcal { O } _ { \phi ^ { 0 } } = \tilde { \psi } _ { f } \zeta ^ { I }, \qquad \mathcal { O } _ { \phi ^ { \prime } } = \left ( q ^ { * } \right ) ^ { A } ( \sigma ^ { \prime } ) ^ { B } q _ { \mathrm { ~ B } } ^ { I }, \quad r = 1, 2, 3 \qquad ( 6 ) } \\ { \mathrm { ~ o t h e ~ a d. ~ b e}}\end{array}
\end{equation}


Note that the first two operators correspond to mass terms for the scalars and fermions
respectively, in the hypermultiplet. The third operator is a triplet with respect to the

\end{document}
